\documentclass[12pt]{article}
\usepackage[utf8]{inputenc}
\usepackage[T1]{fontenc}
\usepackage{hyperref}
\title{Urum Language Resources for K-12 Curriculum Development}
\author{Compiled by Claude Sonnet 5 (claude-sonnet-5) for Andr\'as Kornai}
\date{}
\begin{document}
\maketitle

\section{Introduction}
Urum (ISO 639-3: uum) is a Turkic (Kipchak) language of the Greeks of the northern Azov region of Ukraine, descended from Crimean Tatar dialects. UNESCO classes it as Definitely Endangered. Speaker figures vary widely between sources: Wikipedia cites about 190,000 (a 2000 figure), while Verbix lists about 45,000 for the Azovian variety. No standard literary language exists. Only general descriptive sources were found; no school materials turned up.

\section{Primary Educational Resources}
None found.

\section{Linguistic and Reference Materials}
\begin{itemize}
\item Wikipedia article on the Urum language. \url{https://en.wikipedia.org/wiki/Urum_language}
\item Wikipedia article on the Urums. \url{https://en.wikipedia.org/wiki/Urums}
\item Verbix page on Azovian Urum. \url{https://docs.verbix.com/Languages/UrumAzovian}
\item ``The Urums of the North Azov Sea Region'' (Cyberleninka). \url{https://cyberleninka.ru/article/n/the-urums-of-the-north-azov-sea-region}
\item ``The status of Urum'' (Hellenisteukontos blog). \url{https://hellenisteukontos.opoudjis.net/the-status-of-urum/}
\item ``The North Azovian Greeks and Their Languages'' (Mariupol Memory Park). \url{https://www.mariupolmemorypark.space/?p=2188}
\end{itemize}

\section{Religious and Community Resources}
No community organizations or scripture were identified in the search results.

\section{Digital and Audio Resources}
None found.

\section{Curriculum Development Context}
There is little written literature and no standard orthography. Greek, Ukrainian and Russian influence is strong, and the language is used as a first language by adults only. A curriculum would first need documentation and standardization work. The war in the Mariupol region further complicates fieldwork.
\end{document}